Primary ADC keeps $q'$ in FP32 and reconstructs stored codes \cite{pq,jhq}:
\begin{equation}
 T_m[z]=\|q'_m-\mathcal C_m[z]\|_2^2,\qquad
 \primarydist(q,x)=\sum_{m=0}^{M-1}T_m[z_m(x)].
 \label{eq:5-primary-lut}
\end{equation}
An eight-bit label normally needs 256 entries. The useful extra structure is inside each Cartesian subcodebook, beyond the PQ subspace partition.

At $K=256$, $D_s=8$, the analytical codebook is
\begin{equation}
 \begin{aligned}
 a&=\sigma\Phi^{-1}(3/4)\simeq0.67449\sigma,\qquad
 \mathcal C_m=\{-a,+a\}^{8},\\
 \primarydist(q,x)&=\|q'\|_2^2+da^2
       -2a\sum_{j=0}^{d-1}s_j(x)q'_j.
 \end{aligned}
 \label{eq:5-sign-codebook}
\end{equation}
$\Phi$ is the standard normal CDF. Each byte stores eight signs, higher bits first. For $a>0$, primary order equals sign-vector inner-product order. Cartesian structure also permits exact four-/eight-way splits; a learned replacement needs compatibility checks.

\subsection{Factorisation without changing the quantiser}

The high nibble selects the first four coordinates and the low nibble selects the last four: $\mathcal C_m[16h+l]=(A_m[h],B_m[l])$. Their tables are
\begin{equation}
 T_m^{\rm hi}[h]=\|q'_{m,\rm hi}-A_m[h]\|_2^2,\qquad
 T_m^{\rm lo}[l]=\|q'_{m,\rm lo}-B_m[l]\|_2^2.
 \label{eq:5-lut-halves}
\end{equation}
Disjoint squared-L2 sums give
\begin{equation}
 T_m[z]=T_m^{\rm hi}[z\!\gg\!4]
        +T_m^{\rm lo}[z\mathbin{\&}15].
 \label{eq:5-lut-factorisation}
\end{equation}
This stores 32 entries instead of 256, without changing real-arithmetic scores, codes or candidates.

The builder takes coordinate slices from codewords $16h$ and $l$. This requires all codewords sharing a nibble to share that slice. \texttt{verify\_split\_lut} rejects incompatible tables.
An arbitrary learned 256-centroid PQ subcodebook need not satisfy this
condition: dividing its code labels into nibbles does not factor its
centroid coordinates.

The same-candidate FP32 audit in \cref{sec:7-3} measures numerical
and ranking differences between table forms, including changes after
identical primary survivor sets.


\subsection{Grouping, construction work and resource cost}

The table-grouping ablation generalises the split to $\glut\in\{1,2,4,8\}$
contiguous coordinate groups whose bit fields match the Cartesian
digits. For eight bits and $D_s=8$, each group has $2^{8/\glut}$ entries;
the total entries and FP32 storage per query are
\begin{equation}
 n_{\rm LUT}(\glut)=\glut 2^{8/\glut},\qquad S_{\rm LUT}(\glut)=4Mn_{\rm LUT}(\glut)\ \text{bytes}.
 \label{eq:5-lut-cost}
\end{equation}
Independent entry construction evaluates $D_s2^{8/\glut}$ coordinate
contributions per subspace. Two groups therefore reduce table storage
eightfold and coordinate work sixteenfold relative to independently
constructing every full-table entry. These are table-level cost ratios,
not predicted application speedups.

\begin{table}[tbp]
\centering
\caption[Cartesian lookup-table trade-offs]{Cartesian table trade-offs for an eight-bit, eight-coordinate
primary subspace. Combining additions are the algebraic minimum for
joining the looked-up values, excluding accumulation across subspaces;
they are not measured instruction counts.}
\label{tab:5-lut-resources}
\small
\begin{tabular}{@{}rrrrrr@{}}
\toprule
$\glut$ & Entries & FP32 bytes & Lookups & Combining adds & Build terms \\
\midrule
1 & 256 & 1,024 & 1 & 0 & 2,048 \\
2 & 32 & 128 & 2 & 1 & 128 \\
4 & 16 & 64 & 4 & 3 & 32 \\
8 & 16 & 64 & 8 & 7 & 16 \\
\bottomrule
\end{tabular}
\end{table}

The table is built in global workspace and staged cooperatively when
its complete shared-memory footprint fits the queried per-block opt-in
limit $S_{\max}$. For candidate-buffer capacity $C_{\rm buf}$, the
ordinary exact scan and the grouped scan respectively require
\begin{equation}
 S_0=8C_{\rm buf}+8+S_{\rm LUT}(\glut),\qquad
 S_{\rm grp}=8C_{\rm buf}+24+S_{\rm LUT}(\glut)
 \quad\text{bytes}.
 \label{eq:5-shared-footprint}
\end{equation}
These formulas concern the evaluated scan without optional fused
refinement or early-exit scratch. The code selects shared lookup exactly
when the corresponding footprint is at most $S_{\max}$.

In the table sweep, $c=1{,}000$, block sizes are 512 or 1,024, and
$C_{\rm buf}=2{,}048$. Full tables span 96--384\,KiB, versus
12--48\,KiB for two groups. NVIDIA specifies a 99\,KiB per-block limit
for compute capability 12.0 \cite{blackwell_tuning}. With this limit,
the full-table footprint exceeds capacity even at $M=96$
($114{,}696$ bytes in the query-centric path), whereas the largest
two-group footprint is $65{,}560$ bytes in the grouped path.
Thus the source-level placement rule sends full tables to global memory
and two-group tables to shared memory in these configurations. This placement is inferred from the launch rule and published limit.
The measurement records do not include a device-attribute query.

Four and eight groups both store 16 entries per subspace but perform
more lookups and combinations. Extraction, lookup dependencies, staging,
synchronisation and residency jointly determine the measured trade-off.
Multiple scan groups may stage the same table independently.

The $16+16$ representation is a nibble LUT for an eight-bit sign code.
Small LUTs have precedents in Fast Scan, Quicker ADC and IVF-RaBitQ \cite{fastscan,quickeradc,ivfrabitq}; JHQ retains its own two-level score. Our contribution is their GPU integration with retention and refinement. Two groups lead in 22/24 measured cells; other margins are below 1\% (\cref{sec:7-3}).